\section[扩张 TFT 与拓扑 Maxwell 理论]{扩张拓扑场论与拓扑 Maxwell 理论（3 月 4 日）}

\srcpage{33}

正如上次开始讨论的，扩张 TFT 告诉我们，不仅要给 4-流形指派复数、给 3-流形指派向量空间\footnote{TFT 的这一特定方面在本课中一般不会出现。}，还要给闭曲面指派 $\C$-线性范畴。给定以 $\Sigma$ 为边界的 3-流形 $Y$，我们得到一个线性函子 $Z(\Sigma)\to Z(\emptyset)=\Vect$。这个函子给「$Y$ 连同对象 $E\in Z(\Sigma)$ 的数据」指派一个向量空间。对偶地，可以交换入射与出射，于是得到线性函子 $Z(Y)\colon\Vect\to Z(\Sigma)$，它给出一个特定对象 $Z(Y)(\C)$。（TODO：我想这里该有个用层的例子？）

上次还讨论了一些例子，最值得注意的是拓扑 Maxwell 理论。在 A 侧——它讲的是 $U_1$-丛空间的拓扑，我们称之为「$\Pic$」——我们考虑定向曲面 $\Sigma$，选一个复结构，然后考虑 $\Sigma$ 上线丛的空间 $\Pic(\Sigma)$。这个空间的同伦型不依赖复结构，并且非典范地分裂为
\[
  \Pic(\Sigma)\cong \Z\times \mathrm{Jac}(\Sigma)\times BU_1,
\]
其中最后一个因子（我们忽略它）来自到 $BU_1$ 的常值映射（给出平凡线丛）。你可以把它想成叠性（stackiness）：$BU_1$ 试图成为 $\mathrm{pt}/U_1$，给我们线丛的额外自同构。但现在不必操心。

空间 $\mathrm{Jac}(\Sigma)\cong\C^g/\Z^{2g}$，假设 $\Sigma$ 连通且亏格为 $g$。这是因为它作为 $H^1(\Sigma;\cO_\Sigma)/H^1(\Sigma;\Z)$ 出现。我们只关心拓扑的东西，所以最终得到一个环面，它也是 Eilenberg–Mac Lane 空间 $K(H^1(\Sigma;\Z),1)$。

我们有了一个空间，现在应当线性化以产生范畴，所以取一个层范畴。正如有人说「取空间上的函数」时，你可以取不同具体类型的函数——这是个泛函分析问题——我们也可以考虑不同种类的层。我们取局部常值层，于是指派给 $\Sigma$ 的范畴是 $\mathrm{Loc}(\Pic(\Sigma))$。我们还没有进入导出世界（这意味着暂时忽略分解式中的 $BU_1$，因为 $BU_1$ 单连通），所以 $\mathrm{Loc}(\Pic\Sigma)$ 分解为
\[
  \mathrm{Loc}(\Pic\Sigma)\simeq \Vect_\Z\otimes \mathrm{Loc}(\mathrm{Jac}(\Sigma)),
\]
其中 $\Vect_\Z$ 指 $\Z$-分次向量空间。但因为 Jacobian 是环面，它上面的局部系统典范地等同于 $\Lambda:=H_1(\Sigma;\Z)$ 的表示；因此 $\mathrm{Loc}(\Pic(\Sigma))$ 等同于 $\Z$-分次 $\Lambda$-表示的范畴。它在卷积下有一个幺半结构，而这不是 TFT 告诉我们的。

现在看 B 侧，它讲的是 $\C^\times$-局部系统的代数几何。还是曲面 $\Sigma$，但我们甚至不需要选复结构就可以考虑 $\mathrm{Loc}_{\C^\times}(\Sigma)=\Hom(\pi_1(\Sigma),\C^\times)/\C^\times$——不过由于 $\C^\times$ 阿贝尔，共轭作用是平凡的。又因为 $H_1(\Sigma)=\pi_1(\Sigma)^{ab}$，这个局部系统空间等同于
\[
  \Hom(H_1(\Sigma),\C^\times)/\C^\times \cong (H_1(\Sigma;\Z)\otimes_\Z \C^\times)/\C^\times,
\]
其中 $\C^\times$ 同样平凡地作用。我们应当给它附上一个范畴，某个带代数几何味道的东西，例如 $\QCoh$。这究竟是什么意思？$\mathrm{Loc}_{\C^\times}(\Sigma)$ 又像一个环面，但这个额外的 $\C^\times$-作用去不掉，所以层范畴是 $\C[\mathrm{Loc}_{\C^\times}(\Sigma)]$ 在「乘法」幺半结构下的模范畴，再与 $\Rep_{\C^\times}$ 张量。

我们希望这个范畴等价于 A 侧构造的 $\Pic(\Sigma)$ 上局部常值层的范畴。TODO：我漏掉了这一段，但关键要素是把带卷积的代数 $\C[\Lambda]$ 与带逐点乘法的 $\C[T^\vee_\C]$ 等同起来，从而诱导层的等价。

这是 Fourier–Mukai 对偶（Fourier–Mukai duality）的第一个例子，阿贝尔对偶把层的范畴等同起来。不过我们在 A 侧塞进了额外的 $\Z$-分次，在 B 侧塞进了 $\C^\times$-作用。

这是几何 Langlands 对应（geometric Langlands correspondence）的第一个版本，尽管非常简单。它具体是 $G=\mathrm{GL}_1$ 的 Betti 几何 Langlands 等价。

\srcpage{34}

接下来，我们把这与 3-流形的故事联系起来，然后解释数论类比。设 $Y$ 是紧定向 3-流形，$\partial Y\cong\Sigma$。可以把线丛从 $Y$ 限制到 $\Sigma$，这定义映射 $\pi\colon\Pic(Y)\to\Pic(\Sigma)$；然后可以沿 $\pi$ 推前或拉回层。例如，取 $E\in\mathrm{Loc}(\Pic\Sigma)$，我们想得到一个向量空间，做法是拉回到 $\Pic(Y)$ 并取截面（而不是像 $Y$ 闭时那样只取函数）。

在 B 侧，有 $\pi^\vee\colon\mathrm{Loc}_{\C^\times}Y\to\mathrm{Loc}_{\C^\times}(\Sigma)$，所以同样可以拉回 $E^\vee\in\QCoh(\mathrm{Loc}_{\C^\times}(\Sigma))$ 并取截面。Fourier–Mukai 变换包含一个对应 $E\leftrightarrow E^\vee$，诱导同构
\[
  \Gamma(\Pic(Y),\pi^*E) \cong \Gamma(\mathrm{Loc}_{\C^\times}(Y),(\pi^\vee)^*E^\vee).
\]
这一切看起来很花哨，但记住这些空间只是环面，这里的一切其实来自 Poincaré 对偶与特征标对偶的组合。

现在可以允许分歧了。考虑局部系统的子范畴，使限制到 $\Sigma$ 时落在某个特定子集中。这样做的动机之一来自数论：分歧正是数论中一切有趣事情发生的地方。

于是设 $C$ 是有限域 $\F_q$ 上的曲线。设 $x$ 是 $C$ 的 $\F_q$-点；我们要讨论允许在 $x$ 处分歧。设 $K_x$ 是有理函数域在 $x$ 处的完备化，于是 $K_x\cong\F_{q'}((t))$，其中 $\F_q\hookrightarrow\F_{q'}$ 是有限扩张。则 $\Spec K_x$ 是 $C\setminus x$ 周围的形式圆盘（formal disc）。回忆点对应于拓扑中的圆圈，所以纽结（knot）的管状邻域的边界对应于我们这里看的东西。

在 B 侧（Galois 侧），我们限制
\[
  \mathrm{Loc}_1(C\setminus x)\to \mathrm{Loc}_1(D) = \Hom_{cts}\bigl(\Gal(\bar{K}_x/K_x),\Ql^\times\bigr) = \Hom_{cts}\bigl(\Gal^{ab}K_x,\Ql^\times\bigr).
\]
那么 $\Pic(C)$ 呢？我们把它描述为
\[
  \Pic(C) = F^\times\Big\backslash {\prod_{y\in C}}' K_y^\times \Big/ \prod_{y\in C}\cO_y^\times.
\]
现在必须去掉点 $x$。我们取 $C$ 上的限制乘积，但现在在 $C\setminus x$ 上工作，于是得到
\[
  \Pic(C)/(\Z_x = K_x^\times/\cO_x^\times).
\]
这不是什么好空间，就像在 Riemann 曲面 $\Sigma$ 上把 $\Pic(\Sigma)$ 中对应一个点的那份 $\Z$ 杀掉一样。但它对我们有用，所以称它为 $\Pic(C,\widehat{x})/K_x^\times$；这个记号由第三同构定理保证。$\Pic(C,\widehat{x})$ 是 $C$ 上带 $D_x$ 上平凡化的线丛；等价地，是 $C\setminus x$ 上的线丛连同 $D_x^\times$ 上的平凡化。商掉 $K_x^\times$ 就去掉了形式圆盘上的平凡化。

所以我们不必为了研究去心曲线而把曲线挖空：只需知道 $K_x^\times$ 在那里如何作用。去心圆盘 $D_x^\times$ 上的线丛只有一个同构类，但 $D_x^\times$ 上有很多函数，所以这个线丛有很多自同构。于是作为叠，
\[
  \Pic(D_x^\times) \cong \mathrm{pt}/K_x^\times,
\]
\yizhu{原文此处印作 $\Pic(D_x)^\times$，按上下文应为去心圆盘上的线丛 $\Pic(D_x^\times)$。}{A（原文 p.34）}
而研究这个图景等价于研究 $K_x^\times$ 在 $\Pic(C,\widehat{x})$ 上的作用，这感觉远不像抽象废话，而更像经典数论。然而数论学家说不要取无限水平结构。具体地说，$K_x^\times$ 包含 $\cO^\times=\F_{q'}[[t]]$\yizhu{原文把环 $\cO=\F_{q'}[[t]]$ 与其单位群 $\cO^\times$ 混写为 $\cO^\times=\F_{q'}[[t]]$；按上下文，包含关系说的是 $\cO^\times\subset K_x^\times$。译文照原文保留。}{A（原文 p.34）}，而后者包含对应于同余子群（congruence subgroup）$\cO^\times(n)$ 的子环，即形如 $1+t^n(\cdot)$ 的函数，把它想成到水平 $n$ 的平凡化。于是 $K_x^\times$ 在 $\Pic(C,nx)$（大致是 $n$ 阶邻域）上的作用穿过 $K_x^\times/\cO^\times(n)$。

直到最近，处理这件事都没有很好的技术，所以我们来看 $K^\times$ 的光滑表示（smooth representation），意思不是你在 $\C$ 上通常期待的东西，而更偏 $p$ 进：每个向量都被某个 $\cO^{\times N}$（$N\gg 0$）固定。例如，若给被作用的对象离散拓扑，这就强加了连续性，因为它表达的是 $K_x^\times$ 上的拓扑。

TODO：我漏掉了一些东西，即把这里与「曲面作为 3-流形的边界」作比较。但就像那边在体与边界之间搬运层，我们在 $C\setminus x$ 与 $D_x^\times$ 之间推挪层。

这些不同的数据片之间如何对话？在 A 侧，我们有来自 $\Pic(C,\widehat{x})$ 的 $K_x^\times$ 的表示；在 B 侧，我们有 $\mathrm{Loc}_1(\cO^\times)$ 上的层，即 $\cO^\times$ 上函数的模。我们需要 Fourier–Mukai 定理的类比，才能在这两个类似层范畴的东西之间来回。

\srcpage{35}

局部类域论（local class field theory）讲的正是这件事——尽管它通常并不以这种「范畴等价」的方式表述。正如整体类域论通常被表述为极大无分歧阿贝尔 Galois 群与完备化 Picard 群之间的同构（如我们讨论过的），然后用 Pontrjagin 对偶以阿贝尔对偶的方式重新解释；我们会对局部类域论做类似的事，即同构
\[
  \widehat{K_x}^\times \xrightarrow{\ \sim\ } \Gal^{ab}(\bar{K}_x/K_x).
\]
这里左端取射有限完备化。无论如何，仅这个阿贝尔群的同构就告诉我们：$K_x^\times$ 的群环 Pontrjagin 对偶于 $\Hom(\Gal^{ab}K,\Ql^\times)$。后者正是我们所谓的「$\mathrm{Loc}$」。同样，一侧是带卷积的群代数，另一侧是带乘法的群代数。我们的群代数不是在 $\C$ 上，而是在 $\Ql$ 上。这里必须小心拓扑。

现在过渡到模，这告诉我们 $K_x^\times$ 的表示等价于 $\mathrm{Loc}_1$ 上的拟凝聚层。需说明的是，把这件事做精确需要大量工作。但要点是：一侧的 $K_n^\times$ 的表示\yizhu{原文此处写 $K_n^\times$，按上下文应为 $K_x^\times$；译文照原文保留。}{A（原文 p.35）}，被等同于局部系统空间的代数几何，正如拓扑故事中那样。

我们要把上述同构看成把 $K_x^\times$ 的表示与「$\mathrm{GL}_1$ Galois 表示上函数的模」等同起来。我们这里没有 $\Pic(C,\widehat{x})$ 的模空间解释，所以我们取 adele：
\[
  F^\times\Big\backslash {\prod_v}' K_v^\times \Big/ \prod_{v\notin S}\cO_v^\times,
\]
即分歧 idele 类群，它被一个很大的阿贝尔群作用，即 $\prod_{v\in S}K_v^\times$。

下次我们会结束阿贝尔群的部分，把今天在 Riemann 曲面上讨论的对偶理解为谱分解或 Fourier 变换。这将为课程目标的非阿贝尔推广铺路。
